\documentclass[10pt,aspectratio=43]{beamer}
\usepackage{tikz}
\usetheme{Boadilla}
\setbeamertemplate{navigation symbols}{}
\title{Iterative refinement: experiments and bounds}
\subtitle{Generated interaction study}
\author{Research group}
\date{}
\year=2026\month=11\day=5
\newcommand{\fuzzlayerA}[1]{\textrm{#1}}
\newcommand{\fuzzlayerB}[1]{\fuzzlayerA{#1}}
\newcommand{\fuzzlayerC}[1]{\fuzzlayerB{#1}}
\newcommand{\fuzzterm}[1]{\fuzzlayerC{#1}}
\newcommand{\fuzzpair}[2]{\fuzzterm{#1}: #2}
\newcommand{\fuzzunit}{ms}
\let\fuzzoriginalunit\fuzzunit
\newcommand{\fuzzcoefficient}{3}
\newcommand{\fuzzscale}[2][\fuzzcoefficient]{#1\,#2}
\definecolor{fuzzpaint}{RGB}{255,0,255}
\definecolor{fuzzsecond}{RGB}{18,137,67}
\tikzset{fuzzvertex/.style={draw,rounded corners=1pt,fill=blue!10,inner sep=1pt,minimum width=.35cm,minimum height=.25cm}}
\begin{document}
\section{Iterative refinement}
\begin{frame}[t,shrink=5]{Iterative refinement}{Composed dependencies}
\fuzzpair{Experiment}{Timing: 12\fuzzunit; estimate $\fuzzscale{x}+\fuzzscale[2]{y}$.}\par
{\renewcommand{\fuzzunit}{cycles}\let\fuzzterm\textbf\fuzzpair{LocalScope}{Timing: 7\fuzzunit.}\par}
\fuzzpair{RestoredScope}{Timing: 11\fuzzunit.}\par
\begin{itemize}
\item<1-> \begin{itemize}
\item<1-> \[\begin{pmatrix}a&b\\c&d\end{pmatrix}\begin{pmatrix}x\\y\end{pmatrix}\]
\end{itemize}
\par
\begin{exampleblock}{Message passing}
\begin{tikzpicture}\node[draw] {\includegraphics % image argument
[width=20pt,height=10pt,keepaspectratio] % image argument
{fuzz-images/pattern.v2}};\end{tikzpicture}
\end{exampleblock}
\par
\begin{flushleft}
\begin{tikzpicture}\node[draw] {\includegraphics[angle=90,width=18pt]{fuzz-images/photo.jpg}};\end{tikzpicture}
\end{flushleft}
\item<2-> \begin{itemize}
\item \[E=\sum_{i=1}^{n}w_i x_i\]
\end{itemize}
\par
{\small\color{teal}\raggedright
\[\begin{pmatrix}a&b\\c&d\end{pmatrix}\begin{pmatrix}x\\y\end{pmatrix}\]\par}
\par
\begin{description}
\item[Stage 1] \[\begin{aligned}a+b&=c\\ d-e&=f\end{aligned}\]
\end{description}
\item<2-> \begin{description}
\item[Stage 1] Before \includegraphics[draft,width=18pt]{fuzz-images/pattern.v2} after.
\item[Stage 2] {\small\color{black}\raggedleft
\textit{The final diagram} records the order in which information crosses the partition.  The estimate is $\frac{a+b}{1+c}$.\par}
\item[Stage 3] \begin{tabular}{lrr}Method&Cost&Error\\\hline
Fresh0&19&12\\
Cached1&14&9\\
Cached2&10&15\\
Local3&7&20\\
\end{tabular}
\end{description}
\end{itemize}
\par \begin{tikzpicture}[font=\scriptsize]
\fill<2-3>[fuzzpaint] (0,0) rectangle (.8,.3);
\fill<3->[fuzzsecond] (1,0) rectangle (1.8,.3);
\node[anchor=west] at (2,0.15) {Witness};
\end{tikzpicture}
\only<4>{}
\end{frame}
\renewcommand{\fuzzunit}{Hz}
\renewcommand{\fuzzcoefficient}{2}
\renewcommand{\fuzzterm}[1]{{\textit{#1}}}
\section{Constraint propagation}
\begin{frame}[t]{Constraint propagation}
\fuzzpair{MacroStageB}{Timing: 12\fuzzunit, retained: 9\fuzzoriginalunit.}\par
{\renewcommand{\fuzzunit}{local}\fuzzpair{FrameLocal}{Timing: 5\fuzzunit.}\par}
\fuzzpair{FrameRestored}{Timing: 8\fuzzunit; estimate $\fuzzscale{x}$.}\par
\begin{columns}[T]\begin{column}{.46\textwidth}A local update can change the measurements used by several neighboring components.\end{column}\begin{column}{.46\textwidth}\begin{itemize}\item Input state\item Updated state\end{itemize}\end{column}\end{columns}
\end{frame}
\end{document}
